\documentclass{article}
\usepackage[T1]{fontenc}
\usepackage{lmodern}
\usepackage{graphicx}
\usepackage{amsmath,amssymb}
\usepackage[a4paper,margin=2cm]{geometry}
\usepackage[absolute,overlay]{textpos}
\setlength{\TPHorizModule}{1mm}
\setlength{\TPVertModule}{1mm}
\usepackage[indonesian]{babel}
\usepackage{hyperref}
\setlength{\emergencystretch}{2em}
% unit: Halaman 37

\begin{document}

% === Halaman 37 (llm) ===

\section*{Kombinasi Kriptografi dan Steganografi}

\begin{itemize}
    \item Steganografi bukan pengganti kriptografi, tetapi keduanya saling melengkapi.
    \item Keamanan pesan rahasia dapat ditingkatkan dengan menggabungkan kriptografi dan steganografi.
    \item Mula-mula pesan dienkripsi dengan algoritma I kriptografi, selanjutnya pesan terenkripsi disembunyikan di dalam media lain (citra, video, audio, dll).
\end{itemize}

\end{document}
